\paragraph{Copula sống sót:}
Trong khi một Copula thông thường kết nối các hàm phân phối tích lũy, Copula sống sót ($\hat{C}$) kết nối các hàm sống sót biên để tạo thành hàm sống sót đồng thời thông qua công thức:
\bea
\bar{F}(x_1, \dots, x_d) = \hat{C}(\bar{F}_1(x_1), \dots, \bar{F}_d(x_d))
\eea
Về mặt toán học, nếu vectơ $U$ có phân phối Copula $C$, thì Copula sống sót của $C$ chính là hàm phân phối của vectơ $1-U$, trong đó $U:=(F_1(X_1), \dots, F_d(X_d))$. Mối liên hệ trong trường hợp hai chiều được biểu diễn như sau:
\bea
\hat{C}(1-u_1, 1-u_2) = 1 - u_1 - u_2 + C(u_1, u_2)
\eea
\paragraph{Tính đối xứng xuyên tâm:}
Một vectơ ngẫu nhiên được gọi là đối xứng xuyên tâm quanh điểm $a$ nếu $X - a = a - X$. Đối với Copula, tính chất này có nghĩa là $U = 1 - U$. Khi một Copula có tính đối xứng xuyên tâm, Copula sống sót của nó sẽ bằng chính nó ($\hat{C} = C$). Ví dụ điển hình là các Copula của phân phối elip (như Gauss và $t$) đều có tính đối xứng xuyên tâm, trong khi các Copula Archimedean như Gumbel và Clayton thì không.

\paragraph{Phân phối có điều kiện của Copula:}
Đặc tính này cho phép tính xác suất để một biến vượt qua một ngưỡng nhất định khi biết giá trị cụ thể của biến kia (trên thang đo phân vị). Trong hai chiều, hàm phân phối có điều kiện của $U_2$ khi biết $U_1 = u_1$ được tính bằng đạo hàm riêng:
\bea
C_{U_2 \mid U_1}(u_2 \mid u_1)
= P(U_2 \leq u_2 \mid U_1 = u_1)
= \lim_{\delta \to 0}
\frac{C(u_1 + \delta, u_2) - C(u_1, u_2)}{\delta}
= \frac{\partial}{\partial u_1} C(u_1, u_2)
\eea
\paragraph{Mật độ Copula:}
Hầu hết các Copula tham số thường dùng đều có mật độ (ngoại trừ các trường hợp phụ thuộc hoàn hảo như tính đồng điệu). Mật độ Copula được xác định bằng đạo hàm bậc $d$ của hàm $C(u_1, \dots, u_d)$ theo tất cả các thành phần. Mật độ này rất quan trọng trong việc ước lượng tham số bằng phương pháp tối đa hóa xác suất. Hình ảnh mật độ giúp trực quan hóa cấu trúc phụ thuộc; ví dụ, mật độ của $t$ Copula tập trung nhiều khối lượng xác suất ở các góc hơn so với Gauss Copula do có hiện tượng phụ thuộc đuôi.

\paragraph{Tính trao đổi được:}
Một Copula được gọi là trao đổi được nếu giá trị của nó không đổi khi ta hoán vị các đổi số, thể hiện qua công thức $C(u_1, \dots, u_d) = C(u_{\pi(1)}, \dots, u_{\pi(d)})$ với mọi hoán vị $\pi$. Hầu hết các Copula Archimedean chuẩn là trao đổi được, trong khi Gauss và $t$ Copula chỉ trao đổi được khi ma trận tương quan có các hệ số tương quan cặp bằng nhau.

\paragraph{Đặc tính co giãn của Copula giá trị cực trị:}
Các Copula giá trị cực trị có một đặc tính tỷ lệ đặc biệt, thỏa mãn hệ thức $C(u^t) = C^t(u)$ với mọi $t > 0$. Đặc tính này phản ánh cấu trúc phụ thuộc giới hạn của các giá trị cực đại trong các khối dữ liệu.
